% =====================================================================
\chapter{Independencia lineal, bases y rango}
\label{tema:1}
% =====================================================================

Este tema introduce los conceptos que permiten medir cuántas dimensiones utiliza realmente una matriz: independencia lineal, base, dimensión, rango y núcleo. La presentación sigue los §3.1, §3.4 y~§3.5 de Strang~\cite{strang-ila}, las clases~9 a~11 del curso 18.06~\cite{mit1806} y los §2.5 a~§2.7 de~\cite{mml}. Se trabaja siempre en $\R^{n}$, aunque las definiciones valen en cualquier espacio vectorial real.

\section{Combinaciones lineales y subespacios}

\begin{definicion}[Subespacio]
Un subconjunto $W\subseteq\R^{n}$ es un \emph{subespacio vectorial} si contiene el vector nulo y es cerrado para la suma y el producto por escalares: si $u,v\in W$ y $c\in\R$, entonces $u+v\in W$ y $cu\in W$.
\end{definicion}

\begin{definicion}[Subespacio generado]
Dados $v_1,\dots,v_k\in\R^{n}$, una \emph{combinación lineal} de ellos es un vector $c_1v_1+\dots+c_kv_k$ con $c_i\in\R$. El conjunto de todas las combinaciones lineales es un subespacio, que se denota $\gen\{v_1,\dots,v_k\}$ y se llama \emph{subespacio generado} por los $v_i$.
\end{definicion}

\begin{definicion}[Independencia lineal]\label{def:indep}
Los vectores $v_1,\dots,v_k$ son \emph{linealmente independientes} si la única combinación lineal nula es la trivial:
\[
c_1v_1+\dots+c_kv_k=0\ \Longrightarrow\ c_1=\dots=c_k=0.
\]
En caso contrario se dicen \emph{linealmente dependientes}.
\end{definicion}

Si $V=[v_1\mid\cdots\mid v_k]\in\R^{n\times k}$ es la matriz cuyas columnas son los vectores, la condición equivale a que el sistema homogéneo $Vc=0$ tenga solo la solución $c=0$.

\begin{proposicion}\label{prop:dependencia}
Los vectores $v_1,\dots,v_k$ ($k\ge 2$) son linealmente dependientes si y solo si alguno de ellos es combinación lineal de los demás.
\end{proposicion}

\begin{proof}
Si $\sum_i c_iv_i=0$ con algún $c_j\ne 0$, despejando se obtiene $v_j=-\sum_{i\ne j}(c_i/c_j)\,v_i$. Recíprocamente, si $v_j=\sum_{i\ne j}a_iv_i$, la combinación $\sum_{i\ne j}a_iv_i-v_j=0$ es nula y tiene el coeficiente $-1\ne 0$ en $v_j$.
\end{proof}

\section{Bases, dimensión y coordenadas}

\begin{definicion}[Base]
Una \emph{base} de un subespacio $W$ es una familia de vectores de $W$ linealmente independiente que genera $W$.
\end{definicion}

\begin{proposicion}[Coordenadas]\label{prop:coordenadas}
Si $\{u_1,\dots,u_k\}$ es una base de $W$, todo $w\in W$ se escribe de forma única como $w=c_1u_1+\dots+c_ku_k$. Los números $c_1,\dots,c_k$ son las \emph{coordenadas} de $w$ en esa base.
\end{proposicion}

\begin{proof}
La existencia se debe a que la base genera $W$. Si $w=\sum_i c_iu_i=\sum_i c'_iu_i$, restando se obtiene $\sum_i(c_i-c'_i)u_i=0$ y, por independencia, $c_i=c'_i$ para todo $i$.
\end{proof}

\begin{lema}\label{lem:steinitz}
Si $W=\gen\{v_1,\dots,v_m\}$ y $w_1,\dots,w_p\in W$ con $p>m$, entonces $w_1,\dots,w_p$ son linealmente dependientes.
\end{lema}

\begin{proof}
Cada $w_j$ es combinación de los $v_i$: $w_j=\sum_{i=1}^{m}a_{ij}v_i$. Sea $A=(a_{ij})\in\R^{m\times p}$. El sistema homogéneo $Ax=0$ tiene $m$ ecuaciones y $p>m$ incógnitas; la eliminación gaussiana deja al menos $p-m$ variables libres, luego existe $x\ne 0$ con $Ax=0$. Para ese $x$,
\[
\sum_{j=1}^{p}x_jw_j=\sum_{j=1}^{p}x_j\sum_{i=1}^{m}a_{ij}v_i=\sum_{i=1}^{m}\Bigl(\sum_{j=1}^{p}a_{ij}x_j\Bigr)v_i=\sum_{i=1}^{m}(Ax)_i\,v_i=0,
\]
que es una combinación lineal nula no trivial de los $w_j$.
\end{proof}

\begin{teorema}[Dimensión]\label{teo:dimension}
Todas las bases de un subespacio $W$ tienen el mismo número de elementos. Ese número se llama \emph{dimensión} de $W$ y se denota $\dim W$. En particular, $\dim\R^{n}=n$.
\end{teorema}

\begin{proof}
Sean $\{v_1,\dots,v_m\}$ y $\{w_1,\dots,w_p\}$ dos bases de $W$. Si fuera $p>m$, el \cref{lem:steinitz} aplicado a $W=\gen\{v_i\}$ haría dependientes a los $w_j$, lo que contradice que formen base; luego $p\le m$. Intercambiando los papeles, $m\le p$. La base canónica $e_1,\dots,e_n$ muestra que $\dim\R^{n}=n$.
\end{proof}

\begin{ejemplo}[Cambio de coordenadas]\label{ej:coordenadas}
Sean $u_1=(1,1)$ y $u_2=(1,-1)$, que forman una base de $\R^{2}$. Para hallar las coordenadas de $b=(3,1)$ se resuelve $c_1u_1+c_2u_2=b$:
\[
\begin{cases} c_1+c_2=3,\\ c_1-c_2=1,\end{cases}\qquad\Longrightarrow\qquad c_1=2,\quad c_2=1.
\]
El mismo vector tiene coordenadas $(3,1)$ en la base canónica y $(2,1)$ en la base $\{u_1,u_2\}$. En forma matricial, si $U=[u_1\mid u_2]$, las coordenadas son $c=U^{-1}b$.
\end{ejemplo}

\begin{observacion}
Las coordenadas de un embedding dependen de la base en la que se expresan. Un cambio de base invertible del espacio de activaciones de un modelo, compensado en las capas siguientes, no altera sus predicciones; por eso las coordenadas individuales de un embedding no tienen, en general, un significado propio. Esta observación es el punto de partida de los métodos de interpretabilidad del módulo~23, que buscan direcciones del espacio y no coordenadas.
\end{observacion}

\section{Espacio columna, núcleo y rango}

\begin{definicion}[Espacio columna y núcleo]
Sea $A\in\R^{m\times n}$. El \emph{espacio columna} (o imagen) de $A$ es $\Col(A)=\{Ax : x\in\R^{n}\}\subseteq\R^{m}$, el subespacio generado por sus columnas. El \emph{núcleo} de $A$ es $\Nuc(A)=\{x\in\R^{n}: Ax=0\}\subseteq\R^{n}$.
\end{definicion}

\begin{definicion}[Rango]
El \emph{rango} de $A$ es $\rango(A)=\dim\Col(A)$, el número máximo de columnas linealmente independientes de~$A$.
\end{definicion}

\begin{teorema}[Rango de filas igual a rango de columnas]\label{teo:filas-columnas}
Para toda $A\in\R^{m\times n}$, $\rango(A)=\rango(A\T)$: el número máximo de filas independientes coincide con el de columnas independientes.
\end{teorema}

\begin{proof}
Sea $r=\rango(A)$ y sea $C\in\R^{m\times r}$ la matriz formada por $r$ columnas de $A$ que constituyen una base de $\Col(A)$. Cada columna $a_j$ de $A$ es combinación de las columnas de $C$, $a_j=Cr_j$ con $r_j\in\R^{r}$; agrupando, $A=CR$ con $R=[r_1\mid\cdots\mid r_n]\in\R^{r\times n}$. Leída por filas, esta igualdad dice que cada fila de $A$ es combinación lineal de las $r$ filas de $R$. Por el \cref{lem:steinitz}, a lo sumo $r$ filas de $A$ son independientes: $\rango(A\T)\le\rango(A)$. Aplicando el mismo argumento a $A\T$ se obtiene la desigualdad contraria.
\end{proof}

La factorización $A=CR$ de la demostración es la que Strang usa como punto de partida en~\cite[§1.4]{strang-ila}.

\begin{teorema}[Rango y nulidad]\label{teo:rango-nulidad}
Para toda $A\in\R^{m\times n}$,
\[
\rango(A)+\dim\Nuc(A)=n.
\]
\end{teorema}

\begin{proof}
Sea $\{k_1,\dots,k_p\}$ una base de $\Nuc(A)$ y complétese con $w_1,\dots,w_q$ hasta una base de $\R^{n}$, de modo que $p+q=n$. Se demuestra que $\{Aw_1,\dots,Aw_q\}$ es una base de $\Col(A)$, lo que da $\rango(A)=q=n-p$.

\emph{Generan.} Todo $x\in\R^{n}$ se escribe $x=\sum_i a_ik_i+\sum_j b_jw_j$, y como $Ak_i=0$, $Ax=\sum_j b_j\,Aw_j$.

\emph{Son independientes.} Si $\sum_j b_jAw_j=0$, entonces $A\bigl(\sum_j b_jw_j\bigr)=0$, luego $\sum_j b_jw_j\in\Nuc(A)$ y se escribe $\sum_j b_jw_j=\sum_i a_ik_i$. Como $\{k_i\}\cup\{w_j\}$ es una base de $\R^{n}$, todos los coeficientes son nulos; en particular, $b_j=0$.
\end{proof}

\begin{ejemplo}\label{ej:M}
Sea
\[
M=\begin{pmatrix}1&2&3\\2&4&6\\1&0&1\end{pmatrix}.
\]
La segunda fila es el doble de la primera, y las filas primera y tercera no son proporcionales; luego $\rango(M)=2$. Por el \cref{teo:rango-nulidad}, $\dim\Nuc(M)=3-2=1$. El vector $(1,1,-1)$ cumple $M(1,1,-1)=(1+2-3,\,2+4-6,\,1+0-1)=0$, luego $\Nuc(M)=\gen\{(1,1,-1)\}$. La matriz aplasta $\R^{3}$ sobre el plano $\Col(M)$ y anula la dirección $(1,1,-1)$.
\end{ejemplo}

\section{Matrices de rango 1 y rango de un producto}

\begin{proposicion}\label{prop:rango1}
Si $u\in\R^{m}$ y $v\in\R^{n}$ son no nulos, la matriz $uv\T\in\R^{m\times n}$ tiene rango~1. Recíprocamente, toda matriz de rango~1 es de esta forma.
\end{proposicion}

\begin{proof}
La columna $j$ de $uv\T$ es $v_j\,u$: todas las columnas son múltiplos de $u$, y alguna es no nula porque algún $v_j\ne 0$. Luego $\Col(uv\T)=\gen\{u\}$ y el rango es~1. Si $\rango(A)=1$ y $u$ genera $\Col(A)$, cada columna es $a_j=v_ju$ para cierto $v_j\in\R$, y $A=uv\T$ con $v=(v_1,\dots,v_n)$.
\end{proof}

\begin{proposicion}[Rango de un producto]\label{prop:rango-producto}
Si $B\in\R^{m\times r}$ y $A\in\R^{r\times n}$, entonces $\rango(BA)\le\min\{\rango(A),\rango(B)\}\le r$.
\end{proposicion}

La demostración se propone en el \cref{eje:1-producto}. De la factorización $A=CR$ del \cref{teo:filas-columnas} se deduce además que toda matriz de rango $r$ es suma de $r$ matrices de rango~1, $A=\sum_{i=1}^{r}c_i\,\rho_i\T$, donde $c_i$ es la columna $i$ de $C$ y $\rho_i\T$ la fila $i$ de $R$. El tema~4 obtiene una descomposición de este tipo con piezas ortogonales y ordenadas por importancia.

\begin{observacion}[Matrices de bajo rango en un LLM]\label{obs:lora-rango}
El ajuste fino LoRA~\cite{hu-lora}, que se estudia en el módulo~20, mantiene congelada una matriz de pesos $W_0\in\R^{d_{\text{out}}\times d_{\text{in}}}$ y aprende una corrección $\Delta W=BA$ con $B\in\R^{d_{\text{out}}\times r}$ y $A\in\R^{r\times d_{\text{in}}}$. Por la \cref{prop:rango-producto}, $\rango(\Delta W)\le r$, y el número de parámetros entrenables es $r(d_{\text{in}}+d_{\text{out}})$ en lugar de $d_{\text{in}}d_{\text{out}}$. Con $d_{\text{in}}=d_{\text{out}}=4096$ y $r=8$ se entrenan $65\,536$ parámetros, el $0{,}39\,\%$ de los $16\,777\,216$ de la matriz completa.
\end{observacion}

\begin{lecturas}
\item Strang~\cite{strang-ila}, §3.4 «Independence, Basis, and Dimension», §3.5 «Dimensions of the Four Subspaces» y apéndice~1 «The Ranks of $AB$ and $A+B$».
\item MIT 18.06~\cite{mit1806}, clases~9 «Independence, basis, and dimension», 10 «The four fundamental subspaces» y 11 «Matrix spaces; rank 1; small world graphs».
\item Deisenroth, Faisal y Ong~\cite{mml}, §2.5 «Linear Independence», §2.6 «Basis and Rank» y §2.7.3 «Image and Kernel» (teorema~2.24).
\end{lecturas}

\section*{Ejercicios del tema 1}
\addcontentsline{toc}{section}{Ejercicios del tema 1}

\begin{ejercicio}{1}{eje:1-vf}
Justifique si son verdaderas o falsas: (a) tres vectores de $\R^{2}$ pueden ser linealmente independientes; (b) si las $n$ columnas de $A\in\R^{m\times n}$ son independientes, entonces $n\le m$; (c) una matriz de $5\times 3$ puede tener rango~4.
\end{ejercicio}

\begin{ejercicio}{2}{eje:1-rango1}
Sean $u=(1,2,3)$ y $v=(1,-1)$. Escriba la matriz $uv\T$, calcule su rango y dé una base de su espacio columna y de su núcleo.
\end{ejercicio}

\begin{ejercicio}{2}{eje:1-mml}
(a) Decida si $x_1=(2,-1,3)$, $x_2=(1,1,-2)$ y $x_3=(3,-3,8)$ son linealmente independientes; si no lo son, escriba una relación de dependencia. (b) Escriba $y=(1,-2,5)$ como combinación lineal de $z_1=(1,1,1)$, $z_2=(1,2,3)$ y $z_3=(2,-1,1)$, y explique por qué los coeficientes son únicos.
\fuente{adaptado de Deisenroth, Faisal y Ong~\cite{mml}, cap.~2, ejercicios~2.10(a) y~2.11.}
\end{ejercicio}

\begin{ejercicio}{2}{eje:1-u4}
Halle el mayor número de vectores linealmente independientes entre $v_1=(1,-1,0)$, $v_2=(1,0,-1)$, $v_3=(0,1,-1)$ y $v_4=(1,1,-2)$. Escriba las relaciones de dependencia, dé una base del subespacio que generan y obtenga una ecuación que cumplan todos sus vectores.
\fuente{adaptado de MIT 18.06SC~\cite{mit1806sc}, sesión «Independence, Basis and Dimension», problema~9.1 (vectores de $\R^{3}$ en lugar de $\R^{4}$).}
\end{ejercicio}

\begin{ejercicio}{2}{eje:1-nucleo}
Sea
\[
A=\begin{pmatrix}1&2&0&1\\2&4&1&3\\3&6&1&4\end{pmatrix}.
\]
Calcule $\rango(A)$, una base de $\Col(A)$ y una base de $\Nuc(A)$, y compruebe el \cref{teo:rango-nulidad}.
\end{ejercicio}

\begin{ejercicio}{3}{eje:1-producto}
Demuestre la \cref{prop:rango-producto}. Deduzca que una corrección LoRA $\Delta W=BA$ con $B\in\R^{d\times r}$ no puede igualar una matriz $\Delta W^{\star}\in\R^{d\times d}$ de rango mayor que $r$.
\end{ejercicio}
